\documentclass[ti,nofooter]{own-haw-page}
\usepackage[german]{babel}
\usepackage{color}
\usepackage{enumitem}
\usepackage{nameref}
\usepackage{url}
\usepackage{seqsplit}
\usepackage[utf8]{inputenc}
\usepackage[hidelinks]{hyperref}
\usepackage[many]{tcolorbox}
\usepackage{tabto}

\newcommand{\raus}[1]{}
\newcommand{\aufg}{\marginpar{{\em Protokoll!}}}
\newcommand{\back}{\ensuremath{\backslash}}
\newcommand{\pfeil}{\ensuremath{\rightarrow}}
\newcommand{\platzlinie}[1]{%
~\\[#1]
~\\
%$\underline{\hfill}$~\\
\hrule~\\
}

\definecolor{lightblue}{rgb}{0.8, 0.9 ,1}

\newcommand{\arrow}{{\textgreater{}\textgreater} }

% Tilde%%%%%
\newcommand{\addtildealt}{\textasciitilde{}}
% schicker, funktioniert noch nicht:
%\newcommand{\addtilde}{{\Large\textasciitilde{}}}

\newcommand{\addtildealts}{{\fontfamily{ptm}\selectfont\textasciitilde}{}}

\newcommand{\addtilde}{\raisebox{-0.3em}\textasciitilde{}}
% Abstand
\newcommand{\addspace}[1]{\vspace{#1}}

% Codezeile mit $ beginnend
\newcommand{\command}[1]{%
  \vspace{10pt}
  {\ttfamily \fcolorbox{black}{lightblue}{
      \parbox{\textwidth} {
        \vspace{-10pt}
        \def\labelitemi{\$}
        \begin{itemize}[leftmargin=3.5ex, itemsep=-1ex, rightmargin=4ex]
          #1
          \vspace{-10pt}
      \end{itemize}}
  }}\vspace{10pt}}

%%%%%%%%%%%%%%%%%%%%%%%%%%%%%%%%%%%%%%%%%%
% Neues Design (tcolorbox)
\tcbset{settings/.style={
    fontupper=\ttfamily,
    breakable,
    enhanced, 
    colback=lightblue, 
    arc=0mm, 
    boxrule=0.5pt,
    enlarge bottom by=0.5ex,
    enlarge top by=1ex,  
    left=1.5ex,
    right=0ex,
    top=0.4ex,
    bottom=0.4ex}}

% Nichtnummerierter Code (neu)
\newcommand{\codespace}[1]{
  \begin{tcolorbox}[settings]{}
    #1
  \end{tcolorbox}
  \vspace{10pt}
}

% Codezeile mit $ beginnend (neu, unbenutzt)
\newcommand{\codedollar}[1]{
  \begin{tcolorbox}[settings]{}
    \def\labelitemi{\$}
    \begin{itemize}[leftmargin=3ex, itemsep=-1ex]
      #1
    \end{itemize}
  \end{tcolorbox}
}

% Unbenutzt
\newcommand{\codenumber}[1]{
  \begin{tcolorbox}[settings]{}
    \begin{enumerate}[leftmargin=3ex, label={\arabic*}, itemsep=-1ex]
      #1
    \end{enumerate}
  \end{tcolorbox}
}
%%%%%%%%%%%%%%%%%%%%%%%%%%%%%%%%%%%%%%%%%%

% Nichtnummerierter Code (alt)
\newcommand{\codebox}[1]{%
  \vspace{10pt}
  {\ttfamily \fcolorbox{black}{lightblue}{
      \parbox{\textwidth} {
        \vspace{-10pt}
        \def\labelitemi{}
        \begin{itemize}[leftmargin=1ex, itemsep=-1ex, rightmargin=4ex]
          #1
          \vspace{-10pt}
      \end{itemize}}
  }}\vspace{10pt}}

% Normaler nummerierter Code
\newcommand{\code}[1]{
  \vspace{10pt}
  %{\setlength{\fboxrule}{1pt} % Rahmendicke definieren
  {\ttfamily \fcolorbox{black}{lightblue}{
      \parbox{\textwidth} {
        \vspace{-10pt}
        \begin{enumerate}[leftmargin=3.5ex, label={\arabic*}, itemsep=-1ex, rightmargin=4ex]
          #1
          \vspace{-10pt}
      \end{enumerate}}
  }}\vspace{10pt}}
%}

\sloppy 



\begin{document}
%\date{16.~April~2020}
%\date{24.~April~2020}
\date{22.~November~2021}


%\setcounter{section}{0}
%\setcounter{section}{1}
\setcounter{section}{6}
\newcounter{dummymcpdsv} 
\setcounter{dummymcpdsv}{\value{section}}
\addtocounter{dummymcpdsv}{1}


\begin{page}

\begin{Large}
{\centering \bfseries \"{U}bungsblatt \arabic{dummymcpdsv}\\ zur Veranstaltung  
\glqq RobE -- Einf\"{u}hrung in die Robotik''\\}
\end{Large}
~\\[-0.5cm]


\section{Taskgenerierung und Perzeption}

In der \"{U}bung \arabic{dummymcpdsv} geht es um die Generierung von Tasks und der Perzeption von Objekten mittels 3D Sensordaten. \\


\textbf{Zur Bearbeitung der \"{U}bungsaufgaben:} Bearbeiten Sie auf jeden Fall
\textit{alle} \"{U}bungsaufgaben. Ausgenommen hiervon sind lediglich die mit \glqq
optional'' gekennzeichneten Textstellen. Lesen Sie sich die Aufgaben gut durch.
Sollten Sie eine Aufgabe nicht l\"{o}sen k\"{o}nnen, so beschreiben Sie
zumindest, wie weit Sie gekommen sind und auf welche Weise Sie vorgegangen sind.
%
Aufgaben mit der Randbemerkung \textit{``Protokoll!}'' sind 
%Die Aufgaben sind direkt 
im Protokoll zu beantworten.
Als L\"{o}sung der Aufgaben ist %entweder
%alternativ
(je nach Aufgabe)
Programmcode abzugeben
(kommentiert) oder ein kurzer Text.


\subsection{Vorbereitungsaufgaben}

%\begin{enumerate}
%\item 
%\end{enumerate}

-- keine --

\subsection{Praktikumsaufgaben}

\begin{enumerate}

\item Bearbeiten sie die Schritte \textbf{\nameref{sec:setup}} (Anhang \ref{sec:setup}) und \textbf{\nameref{sec:demo}} (Anhang \ref{sec:demo}).

\item Bearbeiten Sie die Schritte \textbf{\nameref{sec:task}}  (Anhang \ref{sec:task}). F\"{u}gen Sie nach dem zweiten Task einen Weiteren hinzu, bei dem sich der Roboterarm um \textit{0.2 m} nach unten bewegt. Wiederholen Sie diesen Schritt mit einer Verschiebung um \textit{0.4 m} nach oben. Waren alle \textit{tasks} erfolgreich? \aufg\

\item Welche Struktur haben die Nachrichten \textit{Vector3Stamped} und \textit{TwistStamped} des Typs \textit{geometry\_msgs}? \aufg\

\item Bearbeiten Sie die Schritte \textbf{\nameref{sec:pickplace}}  (Anhang \ref{sec:pickplace}). Wie viele unterschiedliche \textit{Pick}-Variationen wurden berechnet und wie viele davon erfolgreich? \"{U}berlegen Sie, ob es sinnvoll ist, mehrere Variationen zu berechnen. \aufg\    
 
\item Was sind die vordefinierten Maße der Objekte \textit{table} und \textit{cylinder}? \aufg\ 

\item Bearbeiten Sie die Schritte \textbf{\nameref{sec:perception}}  (Anhang \ref{sec:perception}). Um was f\"{u}r einen Gegenstand handelt es sich bei dem erfassten Objekt? Welche Vorteile und Nachteile hat es, die Aufl\"{o}sung der \textit{3D Occupancy Grid Map} höher zu stellen? \aufg\

\item Was stellt die gefilterte Point Cloud da? \"{U}berlegen Sie, was für einen Nutzen diese haben k\"{o}nnte.  \aufg\

\end{enumerate}



\appendix

%%%%%%%%%%%%%%%%%%%%%%%%%%%%%%%%%%%%%%%%%%%%%%%%%%%%%
%%%%%%%%%%%%%%%%%%%%%%%%%%%%%%%%%%%%%%%%%%%%%%%%%%%%%
%%%%%%%%%%%%%%%%%%%%%%%%%%%%%%%%%%%%%%%%%%%%%%%%%%%%%

\section{Installation of the task constructor}
\label{sec:setup}

Start by downloading and installing the necessary packages and building the source folder:

\code{
\item echo {\dq source /opt/ros/noetic/setup.bash\dq} >{}> \addtilde/.bashrc
\item echo {\dq export PATH=/usr/lib/ccache:\$PATH\dq} >{}> \addtilde/.bashrc
\item source \addtilde/.bashrc
\item sudo apt-get update
\item sudo apt-get install ros-noetic-moveit \\ros-noetic-panda-moveit-config ros-noetic-moveit-visual-tools
\item mkdir -p moveit\_ws/src
\item cd moveit\_ws/src
\item git clone https://github.com/ros-planning/moveit\_tutorials
\item git clone https://github.com/ros-planning/moveit\_task\_constructor
\item rosdep update
\item rosdep install -{}-from-paths . -{}-ignore-src -{}-rosdistro noetic -y
\item catkin build
\item echo {\dq source \addtilde/moveit\_ws/devel/setup.bash\dq} >""> \addtilde/.bashrc
}

%%%%%%%%%%%%%%%%%%%%%%%%%%%%%%%%%%%%%%%%%%%%%%%%%%%%%
\section{Starting the demonstration}
\label{sec:demo}

This package has three different demonstrations that show how the \textit{task constructor} functions. To start, we need to launch the \textit{RVIZ} environment that includes the \textit{Panda robot arm}. It will start in a predefined initial state. \\

Launch the \textit{RVIZ} environment: 

\code{
\setcounter{enumi}{13}
\item source \addtilde/moveit\_ws/devel/setup.bash
\item roslaunch moveit\_task\_constructor\_demo demo.launch
}

A new \textit{Motion Planning Tasks}-window can be seen on the right. Though, there aren’t any tasks to find yet. Will will change that by adding some. \\

Run the following demonstration: 

\code{
\setcounter{enumi}{15}
\item rosrun moveit\_task\_constructor\_demo cartesian
}

Taking a look at the \textit{RVIZ} environment it can be seen, that tasks have now been added. \\

The \textit{Panda robot arm} will automatically start with execute the tasks one by one. Each of them will be calculated, once they are next. The task manager also discriminated between a successful and a failed inverse kinematics calculation. A ``1'' indicates, that a solution has been found. 


%%%%%%%%%%%%%%%%%%%%%%%%%%%%%%%%%%%%%%%%%%%%%%%%%%%%%
\section{Modifying a motion task list}
\label{sec:task}

In order to edit the demonstration, we have to take a look at its source code:

\code{
\setcounter{enumi}{16}
\item cd \addtilde/moveit\_ws/src/moveit\_task\_constructor/demo/src
\item gedit cartesian.cpp}

Its code structure shows us that each task is within brace brackets. This creates a control variable scope for declarations and definitions within that brace block, that are destroyed after going out of scope. \\

We can add a additional relative motion by simply copying a block, inserting it and editing the motion, that is described by direction vector. Naming it correctly will help identifying it in task list. \\

After saving the file, the \textit{catkin workspace} has to be rebuild: \\

\code{
\setcounter{enumi}{18}
\item catkin build
}

We can now test our modified motion task:

\code{
\setcounter{enumi}{19}
\item source \addtilde/moveit\_ws/devel/setup.bash
\item rosrun moveit\_task\_constructor\_demo cartesian
}


%%%%%%%%%%%%%%%%%%%%%%%%%%%%%%%%%%%%%%%%%%%%%%%%%%%%%
\section{Pick-and-place demonstration with the Task Constructor}
\label{sec:pickplace}

The \textit{MoveIt! Task Constructor} also has a \textit{Pick-and-Place} demonstration, though it is a bit more complex than our last example. \\

Start the demonstration:

\code{
\setcounter{enumi}{21}
\item source \addtilde/moveit\_ws/devel/setup.bash
\item roslaunch moveit\_task\_constructor\_demo demo.launch
\item roslaunch moveit\_task\_constructor\_demo pickplace.launch
}

This demonstration consists of multiple files. \\

The tasks for the \textit{Pick-and-Place} tasks can be viewed in:

\code{
\setcounter{enumi}{24}
\item cd \addtilde/moveit\_ws/src/moveit\_task\_constructor/demo/src
\item gedit pick\_place\_task.cpp
}

The operation itself is described in:

\code{
\setcounter{enumi}{26}
\item gedit pick\_place\_demo.cpp
}

Pose names of predefined poses, object dimensions and other important parameters are stored in the file \textit{panda\_config.yaml}, that can be found here:

\code{
\setcounter{enumi}{27}
\item cd \addtilde/moveit\_ws/src/moveit\_task\_constructor/demo/config
\item gedit panda\_config.yaml
}


%%%%%%%%%%%%%%%%%%%%%%%%%%%%%%%%%%%%%%%%%%%%%%%%%%%%%
\section{Using the Perception Pipeline}
\label{sec:perception}

The \textit{Perception Pipeline} allows the visualization of obstacles that were captured with 3D sensors using a 3D occupancy grid. \textit{MoveIt!} in this regard is using \textit{Octomap} to create that map. \\

Before starting, we will increase the resolution of the 3D occupancy grid. \\

The settings can be found in the \textit{launch}-folder of the \textit{panda\_moveit\_config}-package:

\code{
\setcounter{enumi}{29}
\item cd \addtilde/moveit\_ws/src/panda\_moveit\_config/launch
}

Edit the file and set the resolution to \textit{0.01}:

\code{
\setcounter{enumi}{30}
\item gedit sensor\_manager.launch.xml
}

Launch the demo:

\code{
\setcounter{enumi}{31}
\item rosrun moveit\_tutorials detect\_and\_add\_cylinder\_collision\_object\\\_demo
}

We now can see the 3D occupancy map that has detected an obstacle and set an collision around it. Let's verify that by looking at the sensor data. \\

Go to \textit{Add}, select \textit{By topic} and choose \textit{PointCloud2} of the depth-field camera on the top of the menu. But in order to see the sensor data, we need to disable the occupancy grid. For that, go to \textit{Planning Scene} in the \textit{Displays} window, open \textit{Scene Geometry} and deselect \textit{Show Scene Geometry}. \\

To visualize the filtered point cloud, repeat the previously described steps and  select \textit{PointCloud2} of the topic \textit{/filtered\_cloud} instead. Again, deselect or remove the first point cloud. \\

\textbf{NOTE:} It nothing is displayed, just reselect the previous planning scene and point cloud! \\


%%%%%%%%%%%%%%%%%%%%%%%%%%%%%%%%%%%%%%%%%%%%%%%%%%%%%
%%%%%%%%%%%%%%%%%%%%%%%%%%%%%%%%%%%%%%%%%%%%%%%%%%%%%
%%%%%%%%%%%%%%%%%%%%%%%%%%%%%%%%%%%%%%%%%%%%%%%%%%%%%


\end{page}

\end{document}
